
\chapter{Conclusão}
\label{ch:conclusao}

Neste trabalho, foi desenvolvida e implementada a cinemática do manipulador robótico myCobot~280, desde a formulação matemática dos modelos de cinemática direta e inversa até à demonstração do movimento do centro de preensão da garra para diferentes posições alvo, utilizando o MATLAB em comunicação com o CoppeliaSim.

Os objetivos inicialmente definidos foram cumpridos, tendo sido deduzidos os modelos cinemáticos, realizada a sua implementação em MATLAB e estabelecida a integração com o ambiente de simulação. A validação foi efetuada através de exemplos numéricos e de testes de posicionamento do manipulador.

\section*{Formulação}

A cinemática direta foi desenvolvida através de uma abordagem geométrica, baseada na decomposição do manipulador em dois planos.

No plano lateral, a posição do centro de preensão é determinada a partir das projeções dos elos, utilizando os ângulos absolutos associados às juntas J2, J3 e J4, bem como a inclinação da garra. No plano horizontal, a rotação da junta J1 e o desvio lateral $L_5$ permitem determinar as coordenadas cartesianas $(x,y,z)$ relativamente ao referencial da base.

A cinemática inversa foi formulada recorrendo à mesma decomposição geométrica, mas seguindo a ordem inversa dos cálculos.

Inicialmente, no plano horizontal, são determinadas a orientação da base, $q_1$, e a distância frontal $d$. Posteriormente, no plano lateral, fixa-se uma inclinação para a garra e reduz-se o problema a um manipulador planar de dois graus de liberdade, resolvido através da lei dos cossenos. Esta resolução permite determinar os ângulos das juntas J2 e J3, sendo o ângulo de J4 calculado a partir da orientação pretendida para a garra.

Para uma inclinação fixa, a resolução é analítica, sendo verificadas as condições de alcançabilidade e os limites articulares. Por defeito, é utilizada a configuração com o cotovelo para cima.

Quando não existe uma solução válida para a inclinação preferencial de $90^\circ$, o algoritmo procura inclinações alternativas, realizando posteriormente um refinamento por bisseção para encontrar uma orientação próxima da inicialmente pretendida.

Esta abordagem permitiu alcançar posições que seriam inacessíveis com a garra orientada verticalmente para baixo. No segundo exemplo analisado, correspondente a um alvo situado a aproximadamente $30$ cm do eixo da base, foi necessário ajustar a inclinação da garra para cerca de $82{,}27^\circ$.

Em ambos os exemplos, a aplicação da cinemática direta aos ângulos obtidos pela cinemática inversa permitiu recuperar as posições pretendidas dentro da precisão numérica considerada, confirmando a consistência matemática dos modelos desenvolvidos.

\section*{Implementação e resultados}

Os modelos cinemáticos foram implementados em MATLAB através das funções \texttt{direct\_kin} e \texttt{inverse\_kin}, sendo posteriormente integrados com as funções responsáveis pela execução dos movimentos do manipulador.

A função \texttt{move\_joints} permite executar movimentos suaves no espaço das juntas, enquanto a função \texttt{move\_linear} é utilizada para realizar trajetórias aproximadamente retilíneas no espaço cartesiano.

Para movimentos verticais, a função \texttt{plan\_vertical} permite selecionar uma inclinação fixa da garra que seja válida ao longo de toda a trajetória. Por sua vez, a função \texttt{go\_home} permite regressar à configuração inicial, realizando previamente uma subida de segurança quando necessário.

A comunicação entre o MATLAB e o CoppeliaSim foi estabelecida em modo síncrono, permitindo coordenar o envio dos comandos articulares com a evolução da simulação. As funções de movimento utilizam as leituras das juntas para acompanhar a execução e avaliar a posição do centro de preensão.

Na demonstração realizada no simulador, o manipulador foi deslocado até diferentes peças distribuídas sobre a mesa.

Nos sete alvos alcançados, o erro de posicionamento foi inferior a $6$ mm, com um valor máximo de $5{,}9$ mm. As diferenças articulares máximas registadas variaram entre $0{,}4^\circ$ e $1{,}6^\circ$.

O erro do modelo, calculado através da comparação entre a posição fornecida pela cinemática direta e a posição do centro de preensão obtida no CoppeliaSim, situou-se entre $1{,}2$ mm e $4{,}1$ mm, permanecendo abaixo da tolerância de $5$ mm considerada no diagnóstico.

Estes resultados demonstram uma boa correspondência entre o modelo geométrico implementado e o posicionamento do manipulador no ambiente de simulação.

Os testes de alcance também evidenciaram a importância da inclinação da garra. Para alvos situados a poucos centímetros da superfície da mesa, verificou-se um alcance horizontal aproximado de $27$ a $28$ cm relativamente ao eixo da base, mantendo a garra vertical.

Ao permitir a alteração da inclinação da garra, considerando valores até $\psi=30^\circ$, o alcance aumentou para aproximadamente $36$ a $37$ cm nas condições testadas.

Por outro lado, quatro posições situadas a cerca de $39$ a $41$ cm do eixo da base foram rejeitadas pela cinemática inversa antes da execução do movimento, por não terem sido encontradas configurações válidas segundo as restrições implementadas.

Deste modo, os resultados demonstraram não apenas a capacidade de posicionamento do manipulador, mas também a importância da verificação da alcançabilidade antes da execução dos movimentos.

\section*{Limitações}

Apesar dos resultados obtidos, foram identificadas algumas limitações na abordagem desenvolvida.

Em primeiro lugar, os erros de posicionamento podem ser influenciados pelas diferenças entre o modelo geométrico e o modelo utilizado no simulador, bem como pelos desvios entre os ângulos articulares pretendidos e os efetivamente atingidos.

O maior erro de alcance foi registado na peça rosa, uma das mais afastadas da base, sendo compatível com uma maior dificuldade de posicionamento quando o manipulador apresenta uma configuração mais estendida.

Nestas condições, pequenas diferenças articulares podem provocar desvios mais significativos no centro de preensão. Os efeitos da gravidade e o tempo de estabilização também podem contribuir para estas diferenças, embora não tenha sido possível determinar isoladamente a sua influência.

Outra limitação está relacionada com a estratégia de movimento. A interpolação realizada pela função \texttt{move\_joints} não garante uma trajetória retilínea do centro de preensão, sendo necessário recorrer à função \texttt{move\_linear} quando se pretende seguir uma trajetória cartesiana.

Além disso, a resolução da cinemática inversa considera a junta J5 fixa a $0^\circ$ e privilegia, por defeito, a configuração com o cotovelo para cima. Estas escolhas simplificam a implementação, mas limitam o conjunto de configurações articulares exploradas.

Por fim, os resultados apresentados foram obtidos em ambiente de simulação. Embora o CoppeliaSim permita avaliar os algoritmos em condições controladas, a sua utilização não substitui os testes num manipulador físico.

A identificação e localização dos objetos num sistema real também introduziriam fontes adicionais de erro, associadas ao sistema de visão e à calibração entre os diferentes referenciais.

\section*{Trabalho futuro}

Como trabalho futuro, pretende-se adaptar a implementação desenvolvida para a utilização com o manipulador físico myCobot~280, aproveitando os modelos cinemáticos e as funções de movimento já implementados.

Uma das primeiras etapas consistirá na calibração de uma câmara real, de forma a identificar a posição das peças sobre a mesa e converter corretamente as coordenadas detetadas para o referencial da base do manipulador.

Será também necessário adaptar a interface de comunicação utilizada em MATLAB para permitir o envio de comandos diretamente ao robô físico, substituindo a comunicação com o ambiente de simulação.

Outra etapa importante será a calibração das juntas, incluindo a verificação dos respetivos pontos de referência e a avaliação de possíveis diferenças entre o modelo geométrico e o comportamento real do manipulador.

Deverão ainda ser estudadas estratégias para compensar os efeitos da gravidade, as folgas mecânicas e os erros de posicionamento, particularmente em configurações mais estendidas.

Por fim, será importante reforçar as condições de segurança durante a execução dos movimentos, através da limitação das velocidades e das forças aplicadas, bem como da verificação de possíveis colisões com os objetos e com o ambiente de trabalho.

Estas melhorias permitirão evoluir da validação em simulação para a implementação experimental de um sistema de manipulação robótica, integrando cinemática, planeamento de movimentos e visão artificial.
